% الجزء: sets-functions-relations
% الفصل: sets
% القسم: important-sets

\documentclass[../../../include/open-logic-section]{subfiles}

\begin{document}

\olfileid[ar]{sfr}{set}{imp}
\olsection{بعض المجموعات المهمة}

\begin{ex}
أكثر المجموعات التي سنبحث فيها !!{element}sها أشياء رياضية. ومن
أهمها أربع مجموعات أفردت بأسماء، وهي:
\begin{multline*}
    \Nat = \{{\OLMathNumeral{0}}, {\OLMathNumeral{1}}, {\OLMathNumeral{2}}, {\OLMathNumeral{3}}, \ldots\} \\
    \shoveright{\text{مجموعة الأعداد الطبيعية}}\\
    \shoveleft{\Int = \{\ldots, -{\OLMathNumeral{2}}, -{\OLMathNumeral{1}}, {\OLMathNumeral{0}}, {\OLMathNumeral{1}}, {\OLMathNumeral{2}}, \ldots\}} \\
    \shoveright{\text{مجموعة الأعداد الصحيحة}}\\
    \shoveleft{\Rat = \Setabs{\nicefrac{{\OLId{latin-ordinary}{m}}}{{\OLId{latin-ordinary}{n}}}}{{\OLId{latin-ordinary}{m}}, {\OLId{latin-ordinary}{n}} \in \Int\text{ و }{\OLId{latin-ordinary}{n}} \neq {\OLMathNumeral{0}}}}\\
    \shoveright{\text{مجموعة الأعداد النسبية}}\\
    \shoveleft{\Real = (-\infty, \infty)}\\
    \text{مجموعة الأعداد الحقيقية (المتصل العددي)}
\end{multline*}
وكل واحدة منها \emph{لا نهائية}؛ أي إن عدد !!{element}sها لا يتناهى.
والأعداد الطبيعية في اصطلاح هذا الكتاب تشمل الصفر، كما تبين القائمة.

وكلما انتقلنا في هذه المجموعات إلى التالية ضممنا \emph{أعدادًا أخرى}
إلى ما تقدم. فـ$\Nat \subseteq \Int \subseteq \Rat \subseteq \Real$؛
لأن الطبيعي صحيح، والصحيح نسبي، والنسبي حقيقي. ثم إن
$\Nat \subsetneq \Int \subsetneq
\Rat$؛ فـ$-{\OLMathNumeral{1}}$ صحيح وليس طبيعيًا، و$\nicefrac{{\OLMathNumeral{1}}}{{\OLMathNumeral{2}}}$ نسبي وليس صحيحًا.
وأما $\Rat \subsetneq \Real$ فليس بالمثل في الظهور، إذ يقتضي وجود
حقيقيات ليست نسبية\oliflabeldef{sfr:arith:real:realline}{، وسيأتي الرجوع
إلى ذلك في \olref[arith][real]{realline}}{}.

وقد نحتاج أيضًا إلى مجموعة الأعداد الصحيحة الموجبة $\PosInt = \{{\OLMathNumeral{1}},
{\OLMathNumeral{2}}, {\OLMathNumeral{3}}, \dots\}$، وإلى مجموعة أول عددين طبيعيين وحدهما $\Bin = \{{\OLMathNumeral{0}}, {\OLMathNumeral{1}}\}$.
\end{ex}

\begin{tagblock}{compsci}
\begin{ex}[السلاسل]
ومن الأمثلة مجموعة \emph{السلاسل المنتهية} ${\OLId{latin-looped}{A}}^{*}$ على أبجدية ${\OLId{latin-looped}{A}}$.
والسلسلة على ${\OLId{latin-looped}{A}}$ كل متتالية منتهية من عناصر~${\OLId{latin-looped}{A}}$، ويدخل في ذلك
\emph{السلسلة الفارغة $\Lambda$} ضمن السلاسل على~${\OLId{latin-looped}{A}}$ لكل أبجدية~${\OLId{latin-looped}{A}}$. فمثلًا:
\begin{multline*}
\Bin^*
=\{\Lambda,{\OLMathNumeral{0}},{\OLMathNumeral{1}},{\OLMathNumeral{00}},{\OLMathNumeral{01}},{\OLMathNumeral{10}},{\OLMathNumeral{11}},\\
{\OLMathNumeral{000}},{\OLMathNumeral{001}},{\OLMathNumeral{010}},{\OLMathNumeral{011}},{\OLMathNumeral{100}},{\OLMathNumeral{101}},{\OLMathNumeral{110}},{\OLMathNumeral{111}},{\OLMathNumeral{0000}},\ldots\}.
\end{multline*}
فإذا تألفت السلسلة ${\OLId{latin-ordinary}{x}}={\OLId{latin-ordinary}{x}}_{{\OLMathNumeral{1}}}\ldots {\OLId{latin-ordinary}{x}}_{{\OLId{latin-ordinary}{n}}}\in {\OLId{latin-looped}{A}}^{*}$ من ${\OLId{latin-ordinary}{n}}$ «حروف»
من ${\OLId{latin-looped}{A}}$، كان \emph{طولها}~${\OLId{latin-ordinary}{n}}$، وكتبنا $\len{{\OLId{latin-ordinary}{x}}}={\OLId{latin-ordinary}{n}}$.
\end{ex}
\end{tagblock}

\begin{ex}[المتتاليات اللانهائية]
وللمجموعة ${\OLId{latin-looped}{A}}$ أيضًا مجموعة~${\OLId{latin-looped}{A}}^\omega$، وهي مجموعة المتتاليات اللانهائية
من !!{element}s المجموعة~${\OLId{latin-looped}{A}}$. ومعنى المتتالية اللانهائية
${\OLId{latin-ordinary}{a}}_{\OLMathNumeral{1}}{\OLId{latin-ordinary}{a}}_{\OLMathNumeral{2}}{\OLId{latin-ordinary}{a}}_{\OLMathNumeral{3}}{\OLId{latin-ordinary}{a}}_{\OLMathNumeral{4}}\dots$ قائمة تمتد بلا نهاية في اتجاه واحد، وكل واحد من
الأشياء فيها !!a{element} من~${\OLId{latin-looped}{A}}$.
\end{ex}

\end{document}
